\section{问题重述}
\subsection{问题背景}
无线电干扰源的定位与清除是无线电频谱管理的重要任务。监测设备能够确定干扰源的大范围分布区域，
但其精确定位仍需技术人员进入目标区域。为在危险环境中代替技术人员开展定位和清除工作，
本文对搭载于机器狗的干扰源自动定位清除系统进行数学建模与配套代码实现。

\subsection{问题提出}
干扰源目标分布区域已经确定为一个半径为 1800 米的圆形区域，并已知干扰源、机器狗的各项设定与参数，
需要建模解决以下的问题：

问题 1：根据若干检测点坐标及其对某干扰源的示向度，给出采用交会定位法计算多边形
定位区域直径的算法，并判断以该直径为直径的圆能否覆盖定位区域。

问题 2：已知某全向干扰源在一个检测点处测得的示向度，给出第二个检测点的选择
策略，并给出其候选区域。

问题 3：目标区域内有 10--16 个全向干扰源，具体个数未知。由一只机器狗从原点出发，
制定保证全部干扰源被清除且完成时间尽量短的自动搜索、定位和清除策略，并通过模拟器检验。

问题 4：目标区域同时存在全向和定向干扰源，其他条件同问题 3，给出相应策略与算法并通过模拟器检验。

